\documentclass[12pt]{article}
\usepackage[utf8]{inputenc}
\usepackage{amsmath, amssymb} % Math packages
\usepackage{geometry}
\usepackage{hyperref} % Clickable links
\usepackage{graphicx} % Images

\geometry{a4paper, margin=1in}

\title{Calcolo delle probabilita appunti}
\author{Matteo}
\date{\today}

\begin{document}

\maketitle

\section{Introduction}

la materia si occupa di quantificare l'incertezza

concetto di asseganre una probabilita ad un evento.

3 questioni fondamentali:
\begin{enumerate}
    \item cos'è la probabilità
    \item come si stima la probabilità
    \item quali assiomi verifica la probabilità?
\end{enumerate}

punto 3 quello piu matematico, devi capire anche come li calcola

\textbf{esperimento aleatorio}: esperimento il cui risultato non è certo.(es. lancio un dado)

\textbf{esito}: un esperimento, ha un solo esito. Ogni esito esclude ogni altro possibile esito.(1,2,3,4,5,6.)

\textbf{evento}: sottoinsieme di esiti.(es. pari,dispari.)

\textbf{spazio campionario}: un insieme che contiene tutti gli ipotetici risulati dell'esperimento aleaatorio, rappresentati attraverso un opportuno codice. Lo spazio campionario si indica con omega grande, un suo elemento si chiama \textbf{esito} e si indica con omega piccolo.

esiste \textbf{l'evento certo}: che corrisponde a un'affermazione sempre vera indipendentemente dall'esito.

\textbf{l'evento impossibile}: affermazione sempre falsa.

e \textbf{l'evento elementare}: sottoinsieme con un solo esito, vera per un solo esito.


la probabilità si colloca da qualche parte all'interno dell'intervallo [0,1] con 0=falso e 1=vero.

se vedi P significa probabilità.P[5,6] vuol dire probabilita che esca 5 o 6.

slide 10 definizione di addittività numerabile fatta su infiniti insiemi, il prof preferisce considerare un numero finito di insiemi.

delta di Dirac =(completa) usato per rappresentare probabilita un po meno triviali.




\end{document}
